% SOURCE_AUTHORITY: sga5_fr_workpass.tex, lines 7355--7388 (LF-normalized unit).
% SOURCE_SHA256: 791F4EFFC5E02832D5D77ED03518C8156D6F07E4C8238B03545DB93D883FBB28
\begin{statement}{Lema 4.2.2}
Seja \(X\) um objeto da categoria \(P\). As seguintes afirmações são equivalentes:
\begin{enumerate}[label=\textup{(\roman*)}]
\item \(X\) é adaptado a \(J\).
\item Para todo par \((n,p)\) de inteiros positivos ou nulos com \(p\le n+1\), tem-se
\[
J^pF^{\,n-p+1}(X_n)=0 .
\]
\end{enumerate}
\end{statement}

\begin{prooftext}
\((i)\Rightarrow(ii)\): então
\[
J^pF^{\,n-p+1}(X_n)\subset F^{n+1}(X_n)=0.
\]

\((ii)\Rightarrow(i)\): fixemos inteiros \(p,q\ge0\). Se \(p=q=0\), a inclusão é imediata. Se \(q>n+1\), tem-se \(F^q(X_n)=0\) e a inclusão também é imediata. Podemos, portanto, supor que \(p+q>0\) e \(q\le n+1\). Se \(n<p+q\), \(F^{p+q}(X_n)=0\) por definição. Por outro lado,
\[
J^pF^q(X_n)
\]
está contido em \(J^{\,n+1-q}F^q(X_n)\), e este último é nulo por hipótese. Se \(n\ge p+q\), a imagem de \(J^pF^q(X_n)\) em \(X_{p+q-1}\) pelo morfismo de transição está contida em \(J^pF^q(X_{p+q-1})\), que é igual a zero por hipótese. A inclusão
\[
J^pF^q(X_n)\subset F^{p+q}(X_n)
\]
segue-se disso por definição da filtração de sistema projetivo.

Em particular, de (4.2.2) deduz-se que todo subsistema projetivo de um sistema projetivo adaptado a \(J\) é adaptado a \(J\).
\end{prooftext}

\paragraph{Ejemplos.}
a) Um sistema projetivo \(J\)-ádico é adaptado a \(J\). Daí deduz-se que todo subsistema projetivo de um sistema projetivo \(J\)-ádico é adaptado a \(J\).

b) Se um sistema projetivo \(X\) verifica a condição (ML) e é adaptado a \(J\), seu sistema projetivo de imagens universais também é adaptado a \(J\).
